\subsection{Classification of Admissible Neutral Finite Groups}
  \label{subsec:classification-neutral-groups}

  The structural constraints developed in the previous sections impose strong
  restrictions on the class of finite groups that can realise admissible
  neutral generating sets.

  Three ingredients play a central role in this analysis.
  First, generators must satisfy the character-vanishing condition
  \[
    \chi_\rho(s)=0
  \]
  in a common irreducible representation sector.
  Second, the pairwise character relations impose a rigid Gram geometry for
  the generating set.
  Third, the configuration of generators must be compatible with the isotropy
  constraint on the second-moment tensor introduced in
  Section~\ref{subsec:second-moment}.

  Taken together, these conditions drastically reduce the space of admissible
  structures.
  In particular, many candidate groups fail to realise a generator geometry
  compatible with both the Gram constraints and the isotropy requirement.

  The resulting classification can be summarised as follows.

  \begin{theorem}
    Let $G$ be a finite group admitting a generating set $S$ whose elements are
    neutral in a common irreducible representation sector $\rho$, and whose
    pairwise relations satisfy the Gram constraints described above.

    Then the admissible groups belong to one of the following families:
    \[
      Q_8,\qquad \mathrm{Dic}_n\ (n\ge 3),\qquad 2O,\qquad 2I .
    \]
  \end{theorem}

  The first two families correspond to anisotropic configurations of the
  second-moment tensor.
  These include the quaternion group $Q_8$ and the infinite family of binary
  dihedral groups $\mathrm{Dic}_n$.

  In contrast, the last two groups arise from isotropic configurations in which
  the second-moment tensor is proportional to the identity.
  These highly symmetric arrangements correspond to the binary polyhedral
  groups associated with the octahedral and icosahedral symmetries.

  An important consequence of this analysis is the exclusion of the binary
  tetrahedral group $2T$.
  Although $2T$ is a finite subgroup of $SU(2)$, its character relations do not
  allow a configuration of neutral generators compatible with both the Gram
  constraints and the isotropy condition.
  As a result, $2T$ cannot realise an admissible neutral generating set in the
  sense defined above.

  The classification therefore identifies a restricted set of finite
  structures capable of supporting admissible spectral configurations.

  The detailed verification of the Gram relations, second-moment tensors, and
  character tables for the relevant groups is provided in
  Appendix~\ref{sec:appendix-simulation}.
